\documentclass{article}

% Language setting
% Replace `english' with e.g. `spanish' to change the document language
\usepackage[english]{babel}

\usepackage[letterpaper,top=2cm,bottom=5cm,left=3cm,right=3cm,marginparwidth=1.75cm]{geometry}

% Useful packages
\usepackage{amsmath}
\usepackage{graphicx}
\usepackage[colorlinks=true, allcolors=blue]{hyperref}
\usepackage{amssymb} 
\usepackage{fancyhdr}
\usepackage{lipsum} % Para generar texto de ejemplo
\usepackage{multicol}
\pagestyle{fancy} 

\fancyhf{} % Limpiar encabezado y pie de página por defecto
\setlength{\headheight}{70pt} % Ajuste del tamaño del encabezado (puedes aumentar este valor si es necesario)

% Definición del encabezado

\fancyhead[C]{%
    \begin{minipage}[c]{2.0cm}
        \centering
        \includegraphics[width=2.0cm]{imagenes/uta.png}
    \end{minipage}%
    \hfill
    \begin{minipage}[c]{8cm} 
        \centering
        \small {\textbf{UNIVERSIDAD TÉCNICA DE AMBATO}} \\[0.1cm]
        \scriptsize \textit{FACULTAD DE INGENIERÍA EN SISTEMAS ELECTRÓNICA E INDUSTRIAL} \\[0.1cm]
        \scriptsize \textbf{CARRERA SOFTWARE - METODOLOGÍAS AGILES} \\[0.1cm]
        \scriptsize \textbf{CICLO ACADÉMICO: AGOSTO 2025 - ENERO 2026}
    \end{minipage}%
    \hfill
    \begin{minipage}[c]{2.0cm}
        \centering
        \includegraphics[width=2.0cm]{imagenes/fisei.png}
    \end{minipage}%
}

\begin{document}
\begin{center}
    {\large \textbf{INFORME DE TAREA}} \\    
\end{center}

\section{Portada}
\begin{flushleft}
\begin{tabular}{ l l }
    \textbf{Tema:} & 
 \\
    \textbf{Unidad de Organización Curricular:} & PROFESIONAL \\
    \textbf{Nivel y Paralelo:} & Cuarto - B \\
    \textbf{Alumnos participantes:} & Miranda Tarupi Luis Sebastian \\
    \textbf{Asignatura:} &  \\
    \textbf{Docente:} &  \\
\end{tabular}
\end{flushleft}


\section{Informe de la Tarea}
\subsection{Objetivos}
\subsubsection{General}


\subsubsection{Específicos}
\begin{itemize}
    \item 
\end{itemize}


\subsection{Instrucciones}
\begin{itemize}
   
\end{itemize}

\subsection{Listado de equipos, materiales y recursos}
\textbf{Listado de equipos y materiales generales empleados en la tarea:}
\begin{itemize}
    \item 
\end{itemize}

\textbf{TAC (Tecnologías para el Aprendizaje y Conocimiento) empleados en la tarea:}
\begin{itemize}
    \item[$\checkmark$] Plataformas educativas
    \item[$\square$] Simuladores y laboratorios virtuales
    \item[$\square$] Aplicaciones educativas
    \item[$\square$] Recursos audiovisuales
    \item[$\square$] Gamificación
    \item[$\checkmark$] Inteligencia Artificial
    \item Otros (Especifique): \underline{\hspace{5cm}}
\end{itemize}


\subsection{Desarrollo de la Tarea (Actividades a Desarrollar)}

\subsection{Resultados obtenidos}


%Especifique las TAC's utilizadas en la Tarea. 
\subsection{Habilidades blandas empleadas en la tarea}
\begin{itemize}
    \item[$\square$] Liderazgo
    \item[$\square$] Trabajo en equipo
    \item[$\square$] Comunicación asertiva
    \item[$\square$] La empatía
    \item[$\checkmark$] Pensamiento crítico
    \item[$\square$] Flexibilidad
    \item[$\checkmark$] La resolución de conflictos
    \item[$\square$] Adaptabilidad
    \item[$\checkmark$] Responsabilidad
\end{itemize}



\subsection{Conclusiones}

\subsection{Recomendaciones}

\subsection{Referencias}


\vspace{-3em} %reduce el espacio causado por refname vacío
\renewcommand{\refname}{}
\bibliographystyle{ieeetr}
\bibliography{sample}

\end{document}